\documentclass[12pt,a4paper]{article}
\usepackage[UTF8]{ctex}
\usepackage{geometry}
\usepackage{listings}
\usepackage{xcolor}
\usepackage{hyperref}
\usepackage{amsmath}

\geometry{margin=2.5cm}

\lstset{
    language=Python,
    basicstyle=\ttfamily\small,
    keywordstyle=\color{blue}\bfseries,
    commentstyle=\color{green!60!black},
    stringstyle=\color{red},
    numbers=left,
    numberstyle=\tiny\color{gray},
    stepnumber=1,
    numbersep=5pt,
    frame=single,
    breaklines=true,
    showstringspaces=false,
    tabsize=4
}

\title{Halo 系统代码文档：Optimizers 模块}
\author{代码分析文档}
\date{\today}

\begin{document}

\maketitle

\tableofcontents
\newpage

\section{概述}

\texttt{optimizers} 模块是 Halo 系统的核心协调器，负责：
\begin{itemize}
    \item 从 YAML 配置构建工作流图
    \item 创建和管理多个 Worker 进程（每个 GPU 一个）
    \item 调用调度器生成执行计划
    \item 协调 Worker 执行任务，保证依赖关系
    \item 收集结果并更新 Query 状态
\end{itemize}

\section{Optimizer 类}

\subsection{类定义位置}
\texttt{halo/optimizers/halo\_v.py}

\subsection{类概述}

\texttt{Optimizer} 类是多进程编排器，管理整个工作流的执行生命周期。

\subsection{初始化}

\begin{lstlisting}
def __init__(self, config_path: str, **kwargs):
    self.config = load_config(config_path)
    mp.set_start_method("spawn", force=True)
    
    self.device_cnt = torch.cuda.device_count()
    self.processes: List[mp.Process] = []
    self.cmd_queues: List[mp.Queue] = []
    self.result_queues: List[mp.Queue] = []
    self.dp_threshold = 2
    
    # 构建 OP 图
    self.ops, self.start_ops, self.end_ops, self.models = build_ops_from_config(self.config)
    self._create_workers()
    logger.info("Optimizer initialized")
\end{lstlisting}

\begin{description}
    \item[功能] 初始化优化器，加载配置并创建 Worker 进程
    \item[关键步骤]
    \begin{enumerate}
        \item 加载 YAML 配置
        \item 设置多进程启动方式为 "spawn"（Windows 兼容）
        \item 检测可用 GPU 数量
        \item 为每个 GPU 创建命令队列和结果队列
        \item 解析配置构建 OP 图
        \item 创建 Worker 进程
    \end{enumerate}
\end{description}

\subsection{\_create\_workers 方法}

\begin{lstlisting}
def _create_workers(self) -> None:
    """
    Spawn one worker per visible physical GPU.
    Each worker process is restricted to exactly one physical GPU by
    narrowing CUDA_VISIBLE_DEVICES inside the child process.
    """
    # 为每个 Worker 创建队列
    for _ in range(self.device_cnt):
        self.cmd_queues.append(mp.Queue())
        self.result_queues.append(mp.Queue())
    
    # 解析可见的物理 GPU ID
    phys_ids = _visible_physical_gpu_ids()
    if not phys_ids:
        raise RuntimeError("No visible GPUs.")
    phys_ids = phys_ids[: self.device_cnt]
    
    # 创建进程
    for i, physical_gpu_id in enumerate(phys_ids):
        proc = mp.Process(
            target=worker_process,
            args=(i, physical_gpu_id, self.cmd_queues[i], self.result_queues[i]),
            daemon=False,
        )
        self.processes.append(proc)
        proc.start()
\end{lstlisting}

\begin{description}
    \item[功能] 为每个可见的物理 GPU 创建一个独立的 Worker 进程
    \item[关键点]
    \begin{itemize}
        \item 每个 Worker 有独立的命令队列和结果队列
        \item Worker 进程通过设置 \texttt{CUDA\_VISIBLE\_DEVICES} 绑定到特定 GPU
        \item 使用 \texttt{daemon=False} 确保主进程退出时 Worker 正常结束
    \end{itemize}
\end{description}

\subsection{schedule 方法}

\begin{lstlisting}
def schedule(self, queries: List[Query], strategy: str = "heuristic") -> None:
    """
    Build self.workflows by selected strategy.
    """
    self._optimize_queries(queries)
    self.req_id_map = {q.id: q for q in self.queries}
    
    if strategy == "search":
        self.workflows = schedule_search(...)
    elif strategy == "heuristic":
        self.workflows = schedule_heuristic(...)
    elif strategy == "rr":
        self.workflows = schedule_rr(...)
    elif strategy == "levels":
        self.workflows = schedule_by_levels(...)
    elif strategy == "colocate_parents":
        self.workflows = schedule_colocate_parents(...)
    elif strategy == "dp":
        self.workflows = schedule_dp(...)
    else:
        raise ValueError(f"Unknown schedule strategy: {strategy}")
\end{lstlisting}

\begin{description}
    \item[功能] 根据选定的调度策略生成执行计划
    \item[参数]
    \begin{itemize}
        \item \texttt{queries}：查询列表
        \item \texttt{strategy}：调度策略名称
    \end{itemize}
    \item[执行流程]
    \begin{enumerate}
        \item 优化查询顺序（按优先级和长度）
        \item 创建查询 ID 到 Query 对象的映射
        \item 调用对应的调度器生成 workflows
    \end{enumerate}
\end{description}

\subsubsection{\_optimize\_queries 方法}

\begin{lstlisting}
def _optimize_queries(self, queries: List[Query]) -> None:
    """Sort queries by priority DESC, then by prompt length ASC."""
    self.queries = sorted(queries, key=lambda x: (-x.priority, x.prompt_len))
\end{lstlisting}

\begin{description}
    \item[功能] 对查询进行排序优化
    \item[排序规则]
    \begin{enumerate}
        \item 优先级降序（高优先级先执行）
        \item 提示词长度升序（短查询先执行，可能更快完成）
    \end{enumerate}
\end{description}

\subsection{execute 方法}

这是 Optimizer 的核心方法，负责协调所有 Worker 执行任务。

\subsubsection{方法签名}

\begin{lstlisting}
def execute(
    self,
    queries: List[Query] = None,
    return_queries: bool = False,
    skip_exit: bool = False
):
    """
    Dispatch tasks to workers and collect results.
    
    Added dependency guard:
    - Before sending a task to a worker, verify that for ALL queries in that task,
      outputs from ALL parent ops are already available.
    - If not ready, do not dispatch yet; retry after other tasks complete.
    """
\end{lstlisting}

\subsubsection{执行流程}

\begin{enumerate}
    \item \textbf{初始化执行状态}
    \begin{lstlisting}
    if queries is not None:
        self.schedule(queries, strategy="search")
    
    finish_flags = [False] * self.device_cnt      # Worker 是否完成
    inflight = [False] * self.device_cnt          # Worker 是否有任务在执行
    worker_pointer = [0] * self.device_cnt        # 每个 Worker 的 workflow 指针
    \end{lstlisting}
    
    \item \textbf{定义辅助函数}
    
    \textbf{\_cmd\_transfer}：将调度器生成的 task 转换为 Worker 可执行的格式
    \begin{lstlisting}
    def _cmd_transfer(task: Dict) -> Dict:
        """Translate (op, query_ids) to ExecuteInfo(op=..., query_ids=..., prompts=[...])."""
        if task["command"] == "execute":
            op, query_ids = task["params"][0], task["params"][1]
            prompts = []
            for qid in query_ids:
                prompt = self.req_id_map[qid].prompt
                if isinstance(prompt, list):
                    step = self.req_id_map[qid].step
                    prompt = prompt[step]
                # 拼接父节点的输出作为历史
                history_seqs = [
                    self.req_id_map[qid].op_output.get(inp.id, "")
                    for inp in op.input_ops
                ]
                history = "".join(history_seqs)
                prompts.append(prompt + history)
            exe = ExecuteInfo(op=op, query_ids=query_ids, prompts=prompts)
            task["params"] = (exe,)
        return task
    \end{lstlisting}
    
    \textbf{关键逻辑：依赖关系处理}
    \begin{itemize}
        \item 从 \texttt{Query.op_output} 中提取所有父 OP 的输出
        \item 将这些输出拼接到当前 prompt 后面
        \item 实现多步推理：后续 OP 可以看到前序 OP 的结果
    \end{itemize}
    
    \textbf{\_task\_ready}：检查任务是否满足依赖关系
    \begin{lstlisting}
    def _task_ready(task: Dict) -> bool:
        """
        True if this task can be dispatched:
        - Non-execute commands are always ready.
        - For execute(op, query_ids): all parents' outputs exist for every query.
        """
        if task.get("command") != "execute":
            return True
        op, query_ids = task["params"][0], task["params"][1]
        if not getattr(op, "input_ops", None):
            return True  # 起始节点，无依赖
        parent_ids = [p.id for p in op.input_ops]
        for qid in query_ids:
            q = self.req_id_map[qid]
            for pid in parent_ids:
                if pid not in q.op_output:
                    return False  # 某个父节点输出未就绪
        return True
    \end{lstlisting}
    
    \textbf{\_try\_send}：尝试向 Worker 发送下一个任务
    \begin{lstlisting}
    def _try_send(i: int) -> None:
        """
        Attempt to send the current task for worker i if:
          - worker not finished,
          - no task currently in flight on that worker,
          - task exists and all dependencies are satisfied.
        """
        if finish_flags[i] or inflight[i]:
            return
        
        if worker_pointer[i] >= len(self.workflows[i]):
            # 该 Worker 的所有任务已完成
            finish_flags[i] = True
            if not skip_exit:
                self.cmd_queues[i].put(("exit", ()))
            return
        
        task = self.workflows[i][worker_pointer[i]]
        if _task_ready(task):
            self.cmd_queues[i].put(_cmd_transfer(task))
            inflight[i] = True  # 标记为正在执行
    \end{lstlisting}
    
    \item \textbf{初始派发}
    \begin{lstlisting}
    exe_start = time.perf_counter()
    
    # 尝试向所有 Worker 派发第一个任务
    for i in range(self.device_cnt):
        _try_send(i)
    \end{lstlisting}
    
    \item \textbf{主循环：收集结果并继续派发}
    \begin{lstlisting}
    while not all(finish_flags):
        made_progress = False
        for i in range(self.device_cnt):
            if finish_flags[i] or not inflight[i]:
                continue  # 已完成或空闲（无在途任务）
            
            try:
                message = self.result_queues[i].get(timeout=0.1)
            except queue.Empty:
                continue
            
            # 收到消息 -> Worker i 的任务完成
            inflight[i] = False
            made_progress = True
            
            # 处理结果
            if not isinstance(message, dict):
                logger.warning("Worker %d returned non-dict message", i)
                worker_pointer[i] += 1
                _try_send(i)
                continue
            
            cmd = message.get("command")
            if cmd == "execute":
                result = message.get("result", {})
                op_name = result.get("op_name")
                
                # 更新每个查询的输出
                for rec in result.get("item", []):
                    q = self.req_id_map[rec["id"]]
                    q.op_output[op_name] = rec["output"]
                    q.step += 1
                    q.benchmark[op_name] = rec["benchmark"]
                
                # 更新 OP 的性能指标
                if "benchmark" in result and op_name in self.ops:
                    self.ops[op_name].benchmark.update(result["benchmark"])
            
            elif cmd == "error":
                logger.error("Worker %d error: %s", i, message.get("result"))
            
            # 移动指针，尝试发送下一个任务
            worker_pointer[i] += 1
            _try_send(i)
            
            # 任何任务完成后，其他 Worker 可能被解锁
            # 尝试向所有空闲 Worker 派发任务
            for j in range(self.device_cnt):
                if not finish_flags[j]:
                    _try_send(j)
        
        # 如果没有进展，再次尝试派发（避免死锁）
        if not made_progress:
            for j in range(self.device_cnt):
                if not finish_flags[j]:
                    _try_send(j)
    \end{lstlisting}
    
    \item \textbf{返回结果}
    \begin{lstlisting}
    if return_queries:
        return self.queries
    return time.perf_counter() - exe_start
    \end{lstlisting}
\end{enumerate}

\subsubsection{依赖关系保证机制}

这是 \texttt{execute} 方法的核心特性：

\begin{itemize}
    \item \textbf{依赖检查}：在派发任务前，检查该任务的所有父 OP 是否已完成
    \item \textbf{阻塞处理}：如果依赖未满足，任务不会派发，等待父任务完成
    \item \textbf{自动重试}：每当有任务完成，会重新检查所有等待中的任务，看是否满足依赖
    \item \textbf{死锁避免}：通过定期重试和依赖检查，避免系统死锁
\end{itemize}

\subsection{性能分析}

\subsubsection{print\_latency\_percentiles 方法}

\begin{lstlisting}
def print_latency_percentiles(self):
    """Compute and print global P50/P95 latency across queries."""
    all_latencies = []
    for q in self.req_id_map.values():
        start_time = q.create_time
        points = list(q.benchmark.values())
        if not points:
            continue
        end_time = points[-1][-1]  # 最后一个 benchmark 的结束时间
        all_latencies.append(end_time - start_time)
    
    if not all_latencies:
        return
    
    p50 = float(np.percentile(all_latencies, 50))
    p95 = float(np.percentile(all_latencies, 95))
    print(f"Latency Percentiles: P50={p50:.3f}s, P95={p95:.3f}s")
\end{lstlisting}

\begin{description}
    \item[功能] 计算并打印所有查询的延迟百分位数
    \item[计算方式]
    \begin{itemize}
        \item 对每个查询，计算从创建到最后一个 OP 完成的总时间
        \item 统计所有查询的延迟分布
        \item 输出 P50（中位数）和 P95（95 百分位）
    \end{itemize}
\end{description}

\subsection{资源清理}

\subsubsection{exit 方法}

\begin{lstlisting}
def exit(self):
    """Clean up queues and join all worker processes."""
    for q in self.cmd_queues + self.result_queues:
        q.close()
        q.join_thread()
    for p in self.processes:
        p.join()
    logger.info("Optimizer exited")
\end{lstlisting}

\begin{description}
    \item[功能] 清理资源并等待所有 Worker 进程结束
    \item[步骤]
    \begin{enumerate}
        \item 关闭所有队列
        \item 等待队列线程结束
        \item 等待所有 Worker 进程结束
    \end{enumerate}
\end{description}

\section{worker\_process 函数}

\subsection{功能}

这是每个 Worker 进程的入口点：

\begin{lstlisting}
def worker_process(id, device, cmd_queue, result_queue):
    worker = vLLMWorker(id, device, cmd_queue, result_queue)
    worker.run()
\end{lstlisting}

\begin{description}
    \item[功能] 创建 Worker 实例并启动命令循环
    \item[参数]
    \begin{itemize}
        \item \texttt{id}：Worker ID
        \item \texttt{device}：物理 GPU ID
        \item \texttt{cmd\_queue}：命令队列
        \item \texttt{result\_queue}：结果队列
    \end{itemize}
\end{description}

\section{使用示例}

\subsection{基本用法}

\begin{lstlisting}
from halo.optimizers import Optimizer
from halo.components import Query

# 创建优化器
opt = Optimizer("templates/adv_reason_3.yaml")

# 创建查询
queries = [
    Query(0, "What is Machine Learning?"),
    Query(1, "Explain neural networks."),
]

# 执行
opt.schedule(queries, strategy="heuristic")
queries = opt.execute(return_queries=True)

# 查看结果
for q in queries:
    print(f"Query {q.id} final output: {q.op_output}")

# 性能分析
opt.print_latency_percentiles()

# 清理
opt.exit()
\end{lstlisting}

\section{总结}

\texttt{Optimizer} 类是 Halo 系统的核心协调器，主要职责包括：

\begin{itemize}
    \item \textbf{进程管理}：创建和管理多个 Worker 进程
    \item \textbf{任务调度}：调用调度器生成执行计划
    \item \textbf{依赖保证}：确保任务按依赖关系正确执行
    \item \textbf{结果收集}：收集 Worker 的执行结果并更新 Query 状态
    \item \textbf{性能监控}：提供延迟统计和性能分析
\end{itemize}

通过多进程架构和依赖检查机制，Optimizer 实现了高效、可靠的工作流执行。

\end{document}
